\documentclass[10pt]{article}

% Pacotes extras necessários
\usepackage{amsmath}
\usepackage[lmargin=0.5in, rmargin=0.5in, tmargin=0.5in, bmargin=0.5in, includehead, includefoot]{geometry}
\usepackage{amsfonts}
\usepackage[utf8]{inputenc}
\usepackage[portuguese]{babel}
\usepackage{graphicx}
\usepackage{fancyhdr}
\usepackage{setspace}
\usepackage{listings}
\usepackage{url}
\usepackage{enumitem}
\usepackage{appendix}
\usepackage{subcaption}
\usepackage[dvipsnames]{xcolor}
\usepackage{tikz}
\usepackage{makecell}
\usetikzlibrary{positioning}

% Define estilo para listagens de código
\lstdefinestyle{mystyle}{
    language=Matlab,
    inputencoding=latin1,
    extendedchars=true,
    backgroundcolor=\color{white},
    basicstyle=\ttfamily\footnotesize,
    numbers=left,
    numberstyle=\tiny\color{gray},
    numbersep=5pt,
    tabsize=2,
    breaklines=true,
    keywordstyle=\color{blue},
    commentstyle=\color{green!60!black},
    stringstyle=\color{orange},
    frame=single,
    keepspaces=true,
    showspaces=false,
    showstringspaces=false,
    showtabs=false,
}
\lstset{style=mystyle}

\graphicspath{{./images/}}

\setlength{\parindent}{0pt}
\setstretch{1.5}

\DeclareMathOperator{\sen}{sen}
\DeclareMathOperator{\sinc}{sinc}
\newcommand{\Lap}[1]{\mathcal{L}\left\{#1\right\}}
\newcommand{\bm}[1]{\boldsymbol{#1}}

% Cabeçalho e rodapé em todas as páginas
\pagestyle{fancy}
\fancyhf{}
\fancyhead[L]{Relatório das simulações}
\fancyfoot[C]{\thepage}

% Dados do Grupo
\title{
    Trabalho Nº 04 - LS-MRAC \\
    \large COE603 - Controle Adaptativo
}
\author{
    Caio Cesar Leal Verissimo - 119046624 \\
    Leonardo Soares da Costa Tanaka - 121067652 \\
    Lincoln Rodrigues Proença - 121076407 \\
    Engenharia de Controle e Automação - UFRJ \\
    Rio de Janeiro, Rio de Janeiro, Brasil \\
    Julho de 2025
}
\date{}

\begin{document}

\maketitle

% Geração automática do sumário
\tableofcontents
\newpage

\section{Descrição do Algoritmo}

\subsection{Caso \( n^* = 1 \)}

\subsubsection{Hipóteses Básicas}

Seja a planta dada por:

\[
y = P(s)u, \quad P(s) = k \frac{N_p(s)}{D_p(s)}
\]

Já o modelo de referência:

\[
y_m = M(s)r, \quad M(s) = k_m \frac{N_m(s)}{D_m(s)}
\]

E o erro de rastreamento é dado por:

\[
e_0 = y - y_m
\]

Sem perda de generalidade, pode-se assumir \( k_m = 1 \), de modo que:

\[
M(s) = \frac{N_m(s)}{D_m(s)}
\]

Assumimos:
\begin{itemize}
    \item (1) A ordem da planta \( n \) é conhecida.
    \item (2) O grau relativo \( n^* \) é conhecido.
    \item (3) \( N_p(s) \) é Hurwitz (planta de fase mínima).
    \item (4) O sinal de \( k_p \) é conhecido.
\end{itemize}

\textbf{Objetivo:} Encontrar uma lei de controle \( u(t) \) tal que:
\begin{itemize}
    \item O sistema em malha fechada seja estável;
    \item \( \lim_{t \to \infty} e_0(t) = 0 \), onde \( e_0 = y - y_m \), para quaisquer condições iniciais.
\end{itemize}

\subsubsection{Estrutura do Controlador}

Utiliza-se filtros de variáveis de estado (SVF):
\[
\dot{\omega}_1 = A_f \omega_1 + b_f u, \quad \omega_1 \in \mathbb{R}^{n-1}
\]
\[
\dot{\omega}_2 = A_f \omega_2 + b_f y, \quad \omega_2 \in \mathbb{R}^{n-1}
\]

O vetor de regressores é:
\[
\omega^T = [\omega_1^T \;\; y \;\; \omega_2^T \;\; r], \quad \omega \in \mathbb{R}^{2n}
\]

A lei de controle é dada por:
\[
u = \theta^T \omega, \quad \theta^T = [\theta_1^T \;\; \theta_n \;\; \theta_2^T \;\; \theta_{2n}] \in \mathbb{R}^{2n}
\]

Existe um vetor \( \theta^{*T}= [\theta_1^{*T} \;\; \theta_n^{*} \;\; \theta_2^{*T} \;\; \theta_{2n}^{*}]  \) tal que:
\[
y = P(s)u^* = P(s)\theta^{*T} \omega = M(s)r
\]

Erro de parâmetro:
\[ \tilde{\theta} = \theta - \theta^* \]

Desajuste de controle: 
\[ \tilde{u} = \tilde{\theta}^T \omega \]

\subsubsection{Equação do Erro}

Definindo o erro de parâmetros \( \tilde{\theta} = \theta - \theta^* \) e o erro de controle \( \tilde{u} = \tilde{\theta}^T \omega \), então:

\[
y = M(s) \left[ r + \frac{1}{\theta^*_{2n}} \tilde{\theta}^T \omega \right]
\]

Como \( \frac{1}{\theta^*_{2n}} = k_p \), a equação do erro fica:
\[
e_0 = y - y_m = M(s)\left[r + k_p \tilde{\theta}^T \omega\right] = k_p M(s) \left[\tilde{\theta}^T \omega \right]
\]

Se a lei de controle for generalizada como \( u = \theta^T \omega + v \), temos:
\[
e_0 = k_p M(s) \left[u - \tilde{\theta}^{*T} \omega \right] + \varepsilon
\]

A equação do erro pode ser representada em espaço de estados como:
\[
\dot{e} = A e + k_p B [u - \theta^{*T} \omega], \quad e_0 = C e
\]

Este sistema é uma realização não mínima de \( M(s) \), com \( e \in \mathbb{R}^{3n - 2} \). O termo \( \varepsilon \) aparece como contribuição das condições iniciais e decai exponencialmente.

\subsubsection{Critério de Estabilidade (Lema de Meyer–Kalman–Yakubovich)}

Se \( M(s) \) for estritamente positivo real (SPR), então existe \( P = P^T > 0 \) e \( Q = Q^T > 0 \) tal que:
\[
A^T P + P A = -2Q, \quad P B = C^T
\]

\subsubsection{Projeto via Função de Lyapunov}

Considera-se a função de Lyapunov:

\[
2V(e, \tilde{\theta}) = e^T P e + |k_p| \tilde{\theta}^T \Gamma^{-1} \tilde{\theta}, \quad \Gamma = \Gamma^T > 0
\]

Sua derivada é:

\begin{align*}
2\dot{V} &= \dot{e}^T P e + e^T P \dot{e} + |k_p| \left( \dot{\theta}^T \Gamma^{-1} \tilde{\theta} + \tilde{\theta}^T \Gamma^{-1} \dot{\theta} \right) \\
&= \left(A e + k_p B \tilde{\theta}^T \omega \right)^T P e + e^T P \left(A e + k_p B \tilde{\theta}^T \omega \right) + 2|k_p| \tilde{\theta}^T \Gamma^{-1} \dot{\theta} \\
&= e^T A^T P e + \tilde{\theta}^T \omega B^T k_p P e + e^T P A e + k_p e^T P B \tilde{\theta}^T \omega + 2|k_p| \tilde{\theta}^T \Gamma^{-1} \dot{\theta} \\
&= e^T (A^T P + P A) e + 2 k_p \tilde{\theta}^T \omega B^T P e + 2|k_p| \tilde{\theta}^T \Gamma^{-1} \dot{\theta} \\
&= -2 e^T Q e + 2 k_p \tilde{\theta}^T \omega B^T P e + 2|k_p| \tilde{\theta}^T \Gamma^{-1} \dot{\theta}
\end{align*}

Portanto, temos:

\[
\dot{V} = -e^T Q e + k_p \tilde{\theta}^T \omega B^T P e + |k_p| \tilde{\theta}^T \Gamma^{-1} \dot{\theta}
\]

Reescrevendo com o uso de \( \text{sign}(k_p) \):

\[
\dot{V} = -e^T Q e + |k_p| \tilde{\theta}^T \left[ \text{sign}(k_p) \omega B^T e_0 + \Gamma^{-1} \dot{\theta} \right]
\]

Escolhendo a lei de adaptação:

\[
\dot{\theta} = -\text{sign}(k_p) \Gamma \omega e_0
\]

Obtemos:
\[
\dot{V} = -e^T Q e \leq 0
\]

\textbf{Conclusões:}
\begin{itemize}
    \item \( V(t) \) é limitada, não crescente.
    \item \( e, \tilde{\theta} \in \mathcal{L}_\infty \)
    \item Integrando \( \dot{V} \Rightarrow e \in \mathcal{L}_2 \)
    \item \( r(t) \in \mathcal{L}_{\infty} \Rightarrow \dot{e} \in \mathcal{L}_{\infty}\)
    \item \( e \in \mathcal{L}_2\), \( \dot{e} \in \mathcal{L}_\infty \Rightarrow e \to 0 \)
\end{itemize}

\textbf{Convergência de \( \tilde{\theta} \to 0 \)} requer excitação persistente:
\[
\int_t^{t+T} \omega(\tau) \omega^T(\tau) d\tau \geq \alpha I, \quad \forall t \geq t_0, \; \alpha, T > 0
\]

\subsubsection{Erro Auxiliar}

Para conveniência, definimos o erro auxiliar como:
\[
e_a = \operatorname{sign}(k_p) e_0
\]

Assim, a equação do erro pode ser reescrita como:
\[
e_a = |k_p| M(s) [u - \theta^{*T} \omega]
\]

Sabemos que $M(s)$ é SPR (Strictly Positive Real), portanto:
\[
|k_p| M(s)\text{ é SPR}
\]

O termo desconhecido $|k_p|$ é absorvido pelas matrizes $\{A_m, B_m, C_m\}$. Logo, considerando a representação no espaço de estados:

\[
\begin{cases}
\dot{e} = A_m e + B_m [u - \theta^{*T} \omega] \\
e_a = C_m e
\end{cases}
\]

Onde a realização de $M(s)$ é dada por:
\[
|k_p| M(s) = C_m (sI - A_m)^{-1} B_m
\]

Como o sistema $\{A_m, B_m, C_m\}$ é SPR, então existem matrizes:

\[
\begin{cases}
P = P^T > 0 \\
Q = Q^T > 0
\end{cases}
\quad \text{tais que:} \quad
\begin{cases}
A_m^T P + P A_m = -2Q \\
P B_m = C_m^T
\end{cases}
\]

\subsubsection{Projeto de Lyapunov com $e_a$}

Conjunto de equações do sistema adaptativo:

\[
\begin{cases}
\dot{e} = A_m e + B_m \tilde{\theta}^T \omega \\
e_a = C_m e \\
\dot{\theta} = ?
\end{cases}
\]

Função de Lyapunov candidata:
\[
2V(e, \tilde{\theta}) = e^T P e + \tilde{\theta}^T \Gamma^{-1} \tilde{\theta}, \quad \Gamma = \Gamma^T > 0
\]


Derivada temporal da função de Lyapunov:

\[
\begin{aligned}
2\dot{V} &= \dot{e}^T P e + e^T P \dot{e} + \left(\dot{\tilde{\theta}}^T \Gamma^{-1} \tilde{\theta} + \tilde{\theta}^T \Gamma^{-1} \dot{\tilde{\theta}} \right) \\
&= (A_m e + B_m \tilde{\theta}^T \omega)^T P e + e^T P (A_m e + B_m \tilde{\theta}^T \omega) + 2 \tilde{\theta}^T \Gamma^{-1} \dot{\tilde{\theta}} \\
&= e^T A_m^T P e + \tilde{\theta}^T \omega B_m^T P e + e^T P A_m e + e^T P B_m \tilde{\theta}^T \omega + 2 \tilde{\theta}^T \Gamma^{-1} \dot{\tilde{\theta}} \\
&= e^T (A_m^T P + P A_m) e + 2 \tilde{\theta}^T \omega B_m^T P e + 2 \tilde{\theta}^T \Gamma^{-1} \dot{\tilde{\theta}} \\
&= -2 e^T Q e + 2 \tilde{\theta}^T \omega B_m^T P e + 2 \tilde{\theta}^T \Gamma^{-1} \dot{\tilde{\theta}}
\end{aligned}
\]

Portanto:

\[
\dot{V} = -e^T Q e + \tilde{\theta}^T \left[ \omega e_a + \Gamma^{-1} \dot{\tilde{\theta}}\right]
\]

Como $C_m^T e = e_a$, escolhemos a lei de adaptação como:

\[
\dot{\theta} = -\Gamma \omega e_a
\]

Substituindo, temos:

\[
\dot{V} = -e^T Q e \leq 0
\]

\textbf{Conclusão:} o sistema é estável no sentido de Lyapunov.

\subsubsection{Resumo do Algoritmo}

\begin{table}[htbp]
\centering
\renewcommand{\arraystretch}{1.4}
\begin{tabular}{|>{\bfseries}m{3cm}|m{8cm}|c|}
\hline
\textbf{Subsistema} & \textbf{Equação} & \textbf{Ordem} \\
\hline
Planta & $y = P(s)\, u$ & $n$ \\
\hline
Modelo & $y_m = M(s)\, r$ & $n$ \\
\hline
Erro & $e_a = \operatorname{sign}(k_p)\, (y - y_m)$ & \\
\hline
Controle & $u = \theta^T \omega$ & \\
\hline
$\Lambda$-Filters & $\dot{\omega}_1 = A_f \omega_1 + b_f u$ \newline $\dot{\omega}_2 = A_f \omega_2 + b_f y$ & $n - 1$ \\
\hline
Regressor & $\omega^T = \left[ \omega_1^T \quad y \quad \omega_2^T \quad r \right]$ & \\
\hline
Adaptação & $\dot{\theta} = -\Gamma \omega e_a$ & $2n$ \\
\hline
\end{tabular}
\caption{Resumo dos subsistemas, equações e ordens envolvidas no algoritmo de controle adaptativo.}
\label{tab:resumo_algoritmo1}
\end{table}

\vspace{1em}
\noindent
\textbf{Ordem total do sistema:} \fcolorbox{magenta}{white}{$\displaystyle N = 6n - 2$}

\subsection{Caso $n^* = 2$}

A função de transferência \( M(s) \) não pode ser escolhida como SPR (Strictly Positive Real).
Selecionamos um polinômio \( L(s) \) tal que:
\[
M(s) L(s) \text{ é SPR}
\quad \Rightarrow \quad L(s) = (s + \ell_0)
\quad \text{(grau(L) =  1)}
\]
Temos o seguinte modelo de 2ª ordem:
\[
M(s) = \frac{k_m}{(s+a)(s+b)}, \quad L(s) = s + \ell_0
\]
\[
M(s)L(s) = \frac{k_m(s + \ell_0)}{(s + a)(s + b)}
\]

Condição suficiente para SPR:
\[
a \leq \ell_0 \leq b
\]
com polos e zeros devem ser alternados.

A equação do erro pode ser escrita como:
\[
e_a = \operatorname{sign}(k_p) e_0 = |k_p| ML [L^{-1}u - \theta^{*T}\xi]
\]

Selecionamos o controle virtual:
\[
\zeta = \theta^T \xi \Rightarrow u = L(s)[\theta^T\xi]
\quad \text{com} \quad \xi = L(s)^{-1}[\omega]
\]

Erro como função da diferença paramétrica:

\[
e_a = |k_p| ML \left[ \theta^T \xi - \theta^{*T} \xi \right] = |k_p| ML [\tilde{\theta}^T \xi]
\]

Filtro:
\[
\xi = \frac{1}{s + \ell_0} \omega
\quad \Rightarrow \quad \dot{\xi} + \ell_0 \xi = \omega
\]

Controle implementável:
\[
u = (s + \ell_0)\theta^T \xi = (\dot{\theta}^T \xi + \theta^T \dot{\xi}) + \ell_0 (\theta^T \xi)
\]
\[
\Rightarrow u = \dot{\theta}^T \xi + \theta^T (\dot{\xi} + \ell_0 \xi)
= \theta^T \omega + \dot{\theta}^T \xi
\]
Nenhum diferenciador explícito é necessário!

\subsubsection{Projeto de Lyapunov}

Equação do erro:
\[
e_a = |k_p| ML \tilde{\theta}^T \xi
\]

Representação no espaço de estados:
\[
\begin{cases}
\dot{e} = A_m e + B_m [\tilde{\theta}^T \xi] \\
e_a = C_m e
\end{cases}
\]

Observações:
\begin{itemize}
  \item \(\{A_m, B_m, C_m\}\) é uma realização não-mínima de \( |k_p| M(s)L(s) \)
  \item \( e \in \mathbb{R}^{5n - 2} \)
\end{itemize}

Como o sistema é SPR, então:
\[
\exists \quad
\begin{cases}
P = P^T > 0 \\
Q = Q^T > 0
\end{cases}
\quad \text{tais que:} \quad
\begin{cases}
A_m^T P + P A_m = -2Q \\
P B_m = C_m^T
\end{cases}
\]

Função de Lyapunov:
\[
2V(e, \tilde{\theta}) = e^T P e + \tilde{\theta}^T \Gamma^{-1} \tilde{\theta}, \quad \Gamma = \Gamma^T > 0
\]

Derivada temporal:
\[
\begin{aligned}
2\dot{V} &= \dot{e}^T P e + e^T P \dot{e} + \dot{\tilde{\theta}}^T \Gamma^{-1} \tilde{\theta} + \tilde{\theta}^T \Gamma^{-1} \dot{\theta} \\
&= \left(A_m e + B_m \tilde{\theta}^T \xi \right)^T P e + e^T P \left(A_m e + B_m \tilde{\theta}^T \xi \right) + 2 \tilde{\theta}^T \Gamma^{-1} \dot{\theta} \\
&= e^T (A_m^T P + P A_m) e + 2 \tilde{\theta}^T \xi B_m^T P e + 2 \tilde{\theta}^T \Gamma^{-1} \dot{\theta} \\
&= -2 e^T Q e + 2 \tilde{\theta}^T \xi B_m^T P e + 2 \tilde{\theta}^T \Gamma^{-1} \dot{\tilde{\theta}}
\end{aligned}
\]

Portanto:
\[
\dot{V} = -e^T Q e + \tilde{\theta}^T \Gamma^{-1} [ \Gamma \xi e_a+ \dot{\theta}]
\]

Escolhemos a lei de adaptação:
\[
\dot{\theta} = -\Gamma \xi e_a
\]

Resultado:
\[
\dot{V} = -e^T Q e \leq 0
\]

\subsubsection{Considerações Finais}

\begin{itemize}
  \item O regressor original \( \omega \) foi substituído por \( \xi \)
  \item A dinâmica do filtro \( L^{-1}(s) \) foi incorporada
  \item Como \( \omega \in \mathcal{L}_\infty \) e \( L^{-1}(s) \) é BIBO, então \( \xi \in \mathcal{L}_\infty \)
  \item O controle contém o termo extra \( \dot{\theta}^T \xi \), que tende a zero à medida que \( \theta \to \theta^* \)
  \item Esse termo extra atua apenas durante o transiente de identificação
\end{itemize}

\subsubsection{Resumo do Algoritmo}

\begin{table}[htbp]
\centering
\renewcommand{\arraystretch}{1.4}
\begin{tabular}{|>{\bfseries}m{3cm}|m{8cm}|c|}
\hline
\textbf{Subsistema} & \textbf{Equação} & \textbf{Ordem} \\
\hline
Planta & $y = P(s)\, u$ & $n$ \\
\hline
Modelo & $y_m = M(s)\, r$ & $n$ \\
\hline
Erro & $e_a = \operatorname{sign}(k_p)\, (y - y_m)$ & \\
\hline
Controle & $u = \theta^T \omega + \dot{\theta}^T \xi $ & \\
\hline
SV filters & $\dot{\omega}_1 = A_f \omega_1 + b_f u$ \newline $\dot{\omega}_2 = A_f \omega_2 + b_f y$ & $n - 1$ \\
\hline
Regressor & $\omega^T = \left[ \omega_1^T \quad y \quad \omega_2^T \quad r \right]$ & \\
\hline
$\xi$-Filter & $\dot{\xi} = - \ell_0 \xi + \omega, \quad \ell_0 > 0$ & $2n$ \\
\hline
Adaptação & $\dot{\theta} = -\Gamma \xi e_a$ & $2n$ \\
\hline
\end{tabular}
\caption{Resumo dos subsistemas, equações e ordens envolvidas no algoritmo de controle adaptativo.}
\label{tab:resumo_algoritmo2}
\end{table}

\vspace{1em}
\noindent
\textbf{Ordem total do sistema:} \fcolorbox{magenta}{white}{$\displaystyle N = 8n - 2$}

\section{Resultados das Simulações}

Nesta seção, são apresentados os resultados das simulações realizadas com o sistema adaptativo proposto. A estrutura da simulação foi mantida, variando-se apenas um parâmetro por vez em relação à simulação base. Dessa forma, torna-se possível analisar a sensibilidade e influência de cada parâmetro no desempenho do sistema.

As simulações consideram as seguintes definições de planta e modelo de referência:

\begin{itemize}
    \item \textbf{Planta:} \( P(s) = \dfrac{k_p(s + 2)^2}{s^3} \), com \( k_p = \dfrac{1}{3} \)
    \item \textbf{Modelo de referência:} \( M(s) = \dfrac{1}{s + 1} \)
    \item \textbf{Filtro auxiliar:} \( \frac{1}{\Lambda(s)} = \dfrac{1}{(s + 0.5)(s + 1)(s + 1.5)} \)
\end{itemize}

A configuração de simulação base foi definida da seguinte forma:

\begin{itemize}
    \item Ordem do sistema: $n = 3$
    \item Tempo final da simulação: $t_{\text{final}} = 300$ s
    \item Passo de simulação: $st = 0.01$ s
    \item Condições iniciais: $\theta_0 = \mathbf{0}_{6\times1}$, $y_0 = 0$
    \item $\Gamma = 10I$
    \item Sinal de referência: $r(t) =  3 + 10sqw(0.1 \pi t)$
\end{itemize}

A partir desta configuração base, foram realizadas as seguintes simulações adicionais:

\begin{itemize}
    \item Condições iniciais com $\theta(0) = 0{,}95 \cdot \theta^*$
    \item Condições iniciais com $y_0 = 1$
    \item $\Gamma = 100I$
    \item Sinal de referência com componentes senoidais: $r(t) =  3 + 10sqw(0.1 \pi t) + sin(t) + sin(3t)$
\end{itemize}

Os resultados serão discutidos a seguir, com gráficos de comparação entre a saída do sistema e o modelo, além das estimativas dos parâmetros adaptativos e da evolução dos erros. Cada simulação destacará o impacto da variação considerada sobre o comportamento geral do sistema.


\subsection{Simulação \#1 - Base \( (n^* = 1) \)}

\newpage

\subsection{Simulação \#2 - $\theta(0) = 0.95\theta^{*}$ \( (n^* = 1) \)}

\subsection{Simulação \#3 - $y_0 = 1$ \( (n^* = 1) \)}

\newpage

\subsection{Simulação \#4 - $\Gamma = 100I$ \( (n^* = 1) \)}

\subsection{Simulação \#5 - Sinal de Referência \( (n^* = 1) \)}

\newpage


\subsection{Simulação \#6 -  $\theta(0) = 0.95\theta^{*}$ e Sinal de Referência \( (n^* = 1) \)}

\subsection{Simulação \#7 - Base \( (n^* = 2) \)}

\newpage

\subsection{Simulação \#8 - $\theta(0) = 0.95\theta^{*}$ \( (n^* = 2) \)}

\subsection{Simulação \#9 - $y_0 = 1$ \( (n^* = 2) \)}

\newpage

\subsection{Simulação \#10 - $\Gamma = 100I$ \( (n^* = 2) \)}

\subsection{Simulação \#11 - Sinal de Referência \( (n^* = 2) \)}

\newpage


\subsection{Simulação \#12 -  $\theta(0) = 0.95\theta^{*}$ e Sinal de Referência \( (n^* = 2) \)}


\section{Discussão dos Resultados}

Nesta seção são analisados os comportamentos observados nas simulações do controlador MRAC Direto, com foco em robustez, velocidade de adaptação e sensibilidade a parâmetros.

\textbf{Simulações \#1 e \#6 – Base:} O sistema acompanha a referência $y_m(t)$ com erro $e_a \rightarrow 0$ e convergência suave dos parâmetros $\theta(t)$, validando a estabilidade uniforme global do esquema direto, mesmo sem leis de atualização baseadas em regressão explícita.

\textbf{Simulações \#2 e \#7 – $\theta(0) = 0{,}95 \cdot \theta^*$:} Uma estimativa inicial próxima do valor ideal acelera a convergência dos parâmetros e reduz o tempo de transiente. No entanto, o controlador ainda é capaz de adaptar-se eficientemente a partir de estimativas menos precisas, mostrando resiliência da adaptação direta.

\textbf{Simulações \#3 e \#8 – $y_0 = 10$:} A condição inicial elevada gera um transiente mais acentuado, mas o sistema recupera rapidamente o rastreamento da referência, indicando robustez do MRAC Direto frente a perturbações iniciais significativas.

\textbf{Simulações \#4 e \#9 – $\Gamma = 100I$:} O aumento do ganho de adaptação acelera a resposta do sistema, mas também amplifica a sensibilidade a ruídos e variações rápidas, exigindo um compromisso entre velocidade de adaptação e robustez, comum em esquemas diretos.

\textbf{Simulações \#5 e \#10 – Referência Senoidal:} A referência senoidal fornece maior persistência de excitação, o que favorece a atualização dos parâmetros e melhora a convergência. Mesmo com a maior complexidade da referência, o erro permanece pequeno, evidenciando a eficácia adaptativa do controlador.

\textbf{Conclusão:} As simulações confirmam as principais características do MRAC Direto: estabilidade garantida sob certas condições, rastreamento preciso da referência e adaptação eficaz. O sistema se mostra robusto a diferentes condições iniciais e permite ajustes nos parâmetros de adaptação para melhorar o desempenho, sem comprometer a estabilidade.

\end{document}
